\documentclass[10pt,a4paper]{amsart}
\usepackage{amsmath,amssymb,amsthm,graphicx}
% Retained preamble for the Stacks Project chapter {CH} (shared project preamble).
% The chapter file pinned-original-source.tex only carries \input{{preamble}} on its first line,
% so no preamble text lives inside the chapter itself.  The Stacks Project uses one single shared
% preamble for every chapter: the snapshot stages it three times (sources/stacks-chow/preamble.tex,
% sources/stacks-etale-cohomology/preamble.tex and focused-seeds/AGP/sources/stacks-snapshot-20261011/
% preamble.tex) and those copies are identical in size.  The block below is copied mechanically,
% line for line, from sources/stacks-chow/preamble.tex lines 14-18 and 154-217: the Xy-pic setup and
% the native theorem declarations together with the mathematical macros.  Nothing was added, removed
% or redefined; the wrapper class and the amsmath/amssymb/amsthm/graphicx base packages are supplied
% by the assembler.
% For commutative diagrams we use Xy-pic
% Chapter: stacks-cohomology.tex -- Cohomology of Sheaves.
\usepackage[all]{xy}

% We use 2cell for 2-commutative diagrams.
\xyoption{2cell}

% Theorem environments.
%
\theoremstyle{plain}
\newtheorem{theorem}[subsection]{Theorem}
\newtheorem{proposition}[subsection]{Proposition}
\newtheorem{lemma}[subsection]{Lemma}

\theoremstyle{definition}
\newtheorem{definition}[subsection]{Definition}
\newtheorem{example}[subsection]{Example}
\newtheorem{exercise}[subsection]{Exercise}
\newtheorem{situation}[subsection]{Situation}

\theoremstyle{remark}
\newtheorem{remark}[subsection]{Remark}
\newtheorem{remarks}[subsection]{Remarks}

\numberwithin{equation}{subsection}

% Macros
%
\def\lim{\mathop{\mathrm{lim}}\nolimits}
\def\colim{\mathop{\mathrm{colim}}\nolimits}
\def\Spec{\mathop{\mathrm{Spec}}}
\def\Hom{\mathop{\mathrm{Hom}}\nolimits}
\def\Ext{\mathop{\mathrm{Ext}}\nolimits}
\def\SheafHom{\mathop{\mathcal{H}\!\mathit{om}}\nolimits}
\def\SheafExt{\mathop{\mathcal{E}\!\mathit{xt}}\nolimits}
\def\Sch{\mathit{Sch}}
\def\Mor{\mathop{\mathrm{Mor}}\nolimits}
\def\Ob{\mathop{\mathrm{Ob}}\nolimits}
\def\Sh{\mathop{\mathit{Sh}}\nolimits}
\def\NL{\mathop{N\!L}\nolimits}
\def\CH{\mathop{\mathrm{CH}}\nolimits}
\def\proetale{{pro\text{-}\acute{e}tale}}
\def\etale{{\acute{e}tale}}
\def\QCoh{\mathit{QCoh}}
\def\Ker{\mathop{\mathrm{Ker}}}
\def\Im{\mathop{\mathrm{Im}}}
\def\Coker{\mathop{\mathrm{Coker}}}
\def\Coim{\mathop{\mathrm{Coim}}}

% Boxtimes
%
\DeclareMathSymbol{\boxtimes}{\mathbin}{AMSa}{"02}

%
% Macros for moduli stacks/spaces
%
\def\QCohstack{\mathcal{QC}\!\mathit{oh}}
\def\Cohstack{\mathcal{C}\!\mathit{oh}}
\def\Spacesstack{\mathcal{S}\!\mathit{paces}}
\def\Quotfunctor{\mathrm{Quot}}
\def\Hilbfunctor{\mathrm{Hilb}}
\def\Curvesstack{\mathcal{C}\!\mathit{urves}}
\def\Polarizedstack{\mathcal{P}\!\mathit{olarized}}
\def\Complexesstack{\mathcal{C}\!\mathit{omplexes}}
% \Pic is the operator that assigns to X its picard group, usage \Pic(X)
% \Picardstack_{X/B} denotes the Picard stack of X over B
% \Picardfunctor_{X/B} denotes the Picard functor of X over B
\def\Pic{\mathop{\mathrm{Pic}}\nolimits}
\def\Picardstack{\mathcal{P}\!\mathit{ic}}
\def\Picardfunctor{\mathrm{Pic}}
\def\Deformationcategory{\mathcal{D}\!\mathit{ef}}
\usepackage{cleveref}
\crefname{originallemma}{original lemma}{original lemmas}
\Crefname{originallemma}{Original lemma}{Original lemmas}
\makeatletter
\newlabel{algebra-lemma-directed-colimit-exact}{{00DB}{source}{}{sourcepointer-1}{}}
\newlabel{algebra-lemma-directed-colimit-exact@cref}{{[originallemma][0][]00DB}{[1][0][]source}{}{}{}}
\makeatother
\crefname{originallemma}{original lemma}{original lemmas}
\Crefname{originallemma}{Original lemma}{Original lemmas}
\makeatletter
\newlabel{homology-lemma-product-injectives}{{0137}{source}{}{sourcepointer-2}{}}
\newlabel{homology-lemma-product-injectives@cref}{{[originallemma][0][]0137}{[1][0][]source}{}{}{}}
\makeatother
\crefname{originallemma}{original lemma}{original lemmas}
\Crefname{originallemma}{Original lemma}{Original lemmas}
\makeatletter
\newlabel{lemma-cech-vanish-basis}{{01EW}{source}{}{sourcepointer-3}{}}
\newlabel{lemma-cech-vanish-basis@cref}{{[originallemma][0][]01EW}{[1][0][]source}{}{}{}}
\makeatother
\crefname{originallemma}{original lemma}{original lemmas}
\Crefname{originallemma}{Original lemma}{Original lemmas}
\makeatletter
\newlabel{lemma-injective-trivial-cech}{{01EP}{source}{}{sourcepointer-4}{}}
\newlabel{lemma-injective-trivial-cech@cref}{{[originallemma][0][]01EP}{[1][0][]source}{}{}{}}
\makeatother
\crefname{originallemma}{original lemma}{original lemmas}
\Crefname{originallemma}{Original lemma}{Original lemmas}
\makeatletter
\newlabel{lemma-pushforward-injective-flat}{{02N5}{source}{}{sourcepointer-5}{}}
\newlabel{lemma-pushforward-injective-flat@cref}{{[originallemma][0][]02N5}{[1][0][]source}{}{}{}}
\makeatother
\crefname{originallemma}{original lemma}{original lemmas}
\Crefname{originallemma}{Original lemma}{Original lemmas}
\makeatletter
\newlabel{lemma-quasi-separated-cohomology-colimit}{{01FF}{source}{}{sourcepointer-6}{}}
\newlabel{lemma-quasi-separated-cohomology-colimit@cref}{{[originallemma][0][]01FF}{[1][0][]source}{}{}{}}
\makeatother
\crefname{originallemma}{original lemma}{original lemmas}
\Crefname{originallemma}{Original lemma}{Original lemmas}
\makeatletter
\newlabel{sheaves-lemma-descend-opens}{{0A33}{source}{}{sourcepointer-7}{}}
\newlabel{sheaves-lemma-descend-opens@cref}{{[originallemma][0][]0A33}{[1][0][]source}{}{}{}}
\makeatother
\crefname{originallemma}{original lemma}{original lemmas}
\Crefname{originallemma}{Original lemma}{Original lemmas}
\makeatletter
\newlabel{topology-lemma-descend-opens}{{0A30}{source}{}{sourcepointer-8}{}}
\newlabel{topology-lemma-descend-opens@cref}{{[originallemma][0][]0A30}{[1][0][]source}{}{}{}}
\makeatother
\title{Cohomology of a filtered colimit of pulled back sheaves as the colimit of the cohomologies}
\author{The Stacks Project Authors}
\date{stacks-project:cohomology.tex:local-snapshot-20261011; GNU FDL 1.2 or later}
\begin{document}
\maketitle
\noindent\textit{Editorial note.} All delivered mathematical text is the author's own, copied from the pinned Stacks Project chapter file cohomology.tex; no statement or proof step was rewritten, reordered or removed.\par\medskip
\noindent\textit{Editorial note.} The statement begins 'In the situation discussed above.', referring to the setup of the preceding paragraph of the chapter; the counted span reproduces the statement as written and the setup itself is a navigation pointer carried as a stub.\par\medskip
\noindent\textit{Editorial note.} The counted spans carry 8 navigation stubs and no bibliography citation.\par\medskip
\noindent\textit{Editorial note.} The preamble is the chapter's retained minimal preamble (sources/stacks-cohomology/retained-preamble.tex lines 1-81), a verbatim excerpt of the native preamble.\par\medskip
\noindent\textit{Editorial note.} The label lemma-colimit is defined only in cohomology.tex, so no cross-chapter restatement had to be weighed for this card.\par\medskip
\noindent\textit{Allowed standard prerequisites (editorial).}\begin{itemize}
\item Filtered colimits are exact. Algebra, Lemma algebra-lemma-directed-colimit-exact, invoked at proof lines 3427-3428. Supplies that the colimit of the acyclic Cech complexes is acyclic in positive degrees.
\item Products of injective sheaves are injective. Homology, Lemma homology-lemma-product-injectives, invoked at proof lines 3400-3401. Supplies that the sheaves G\_i built as products of pushforwards are injective.
\item Pushforward along a flat morphism of ringed spaces preserves injectives. Same chapter, Lemma lemma-pushforward-injective-flat, invoked at proof lines 3399-3400. Supplies that the individual pushforwards f\_\{b,*\}I\_k are injective on X\_i.
\item Vanishing criterion by Cech cohomology on a basis. Same chapter, Lemma lemma-cech-vanish-basis, invoked at proof lines 3413-3415. Supplies that vanishing of Cech cohomology for coverings of quasi-compact opens by finitely many quasi-compact opens implies vanishing of cohomology.
\item Open coverings of a quasi-compact open descend to some X\_i. Topology, Lemma topology-lemma-descend-opens, invoked at proof lines 3416-3417. Supplies that such a covering is the inverse image of a covering on some X\_i.
\end{itemize}\medskip
\ifdefined\hypertarget\hypertarget{sourcepointer-1}{}\fi
\noindent\textit{Source reference.} Original lemma 00DB: sources/stacks-chow/tags.txt, source lines 491--491. Pointer outside the counted span, invoked once for context; the load-bearing argument of the counted statement is the author's proof reproduced in full in the proof span.\par\medskip
\ifdefined\hypertarget\hypertarget{sourcepointer-2}{}\fi
\noindent\textit{Source reference.} Original lemma 0137: sources/stacks-chow/tags.txt, source lines 1425--1425. Pointer outside the counted span, invoked once for context; the load-bearing argument of the counted statement is the author's proof reproduced in full in the proof span.\par\medskip
\ifdefined\hypertarget\hypertarget{sourcepointer-3}{}\fi
\noindent\textit{Source reference.} Original lemma 01EW: sources/stacks-cohomology/pinned-original-source.tex, source lines 1696--1696. Pointer outside the counted span, invoked once for context; the load-bearing argument of the counted statement is the author's proof reproduced in full in the proof span.\par\medskip
\ifdefined\hypertarget\hypertarget{sourcepointer-4}{}\fi
\noindent\textit{Source reference.} Original lemma 01EP: sources/stacks-cohomology/pinned-original-source.tex, source lines 1408--1408. Pointer outside the counted span, invoked once for context; the load-bearing argument of the counted statement is the author's proof reproduced in full in the proof span.\par\medskip
\ifdefined\hypertarget\hypertarget{sourcepointer-5}{}\fi
\noindent\textit{Source reference.} Original lemma 02N5: sources/stacks-cohomology/pinned-original-source.tex, source lines 1821--1821. Pointer outside the counted span, invoked once for context; the load-bearing argument of the counted statement is the author's proof reproduced in full in the proof span.\par\medskip
\ifdefined\hypertarget\hypertarget{sourcepointer-6}{}\fi
\noindent\textit{Source reference.} Original lemma 01FF: sources/stacks-cohomology/pinned-original-source.tex, source lines 3194--3194. Pointer outside the counted span, invoked once for context; the load-bearing argument of the counted statement is the author's proof reproduced in full in the proof span.\par\medskip
\ifdefined\hypertarget\hypertarget{sourcepointer-7}{}\fi
\noindent\textit{Source reference.} Original lemma 0A33: sources/stacks-chow/tags.txt, source lines 12726--12726. Pointer outside the counted span, invoked once for context; the load-bearing argument of the counted statement is the author's proof reproduced in full in the proof span.\par\medskip
\ifdefined\hypertarget\hypertarget{sourcepointer-8}{}\fi
\noindent\textit{Source reference.} Original lemma 0A30: sources/stacks-chow/tags.txt, source lines 12723--12723. Pointer outside the counted span, invoked once for context; the load-bearing argument of the counted statement is the author's proof reproduced in full in the proof span.\par\medskip
\begin{lemma}
\label{lemma-colimit}
In the situation discussed above.
Let $i \in \Ob(\mathcal{I})$ and let $U_i \subset X_i$ be quasi-compact open.
Then
$$
\colim_{a : j \to i} H^p(f_a^{-1}(U_i), \mathcal{F}_j) =
H^p(p_i^{-1}(U_i), \mathcal{F})
$$
for all $p \geq 0$. In particular we have
$H^p(X, \mathcal{F}) = \colim H^p(X_i, \mathcal{F}_i)$.
\end{lemma}

\begin{proof}
The case $p = 0$ is Sheaves, Lemma \ref{sheaves-lemma-descend-opens}.

\medskip\noindent
In this paragraph we show that we can find a map of systems
$(\gamma_i) : (\mathcal{F}_i, \varphi_a) \to (\mathcal{G}_i, \psi_a)$
with $\mathcal{G}_i$ an injective abelian sheaf and $\gamma_i$ injective.
For each $i$ we pick an injection $\mathcal{F}_i \to \mathcal{I}_i$
where $\mathcal{I}_i$ is an injective abelian sheaf on $X_i$.
Then we can consider the family of maps
$$
\gamma_i :
\mathcal{F}_i
\longrightarrow
\prod\nolimits_{b : k \to i} f_{b, *}\mathcal{I}_k = \mathcal{G}_i
$$
where the component maps are the maps adjoint to the maps
$f_b^{-1}\mathcal{F}_i \to \mathcal{F}_k \to \mathcal{I}_k$.
For $a : j \to i$ in $\mathcal{I}$ there is a canonical map
$$
\psi_a : f_a^{-1}\mathcal{G}_i \to \mathcal{G}_j
$$
whose components are the canonical maps
$f_b^{-1}f_{a \circ b, *}\mathcal{I}_k \to f_{b, *}\mathcal{I}_k$
for $b : k \to j$. Thus we find an injection
$\{\gamma_i\} : \{\mathcal{F}_i, \varphi_a) \to (\mathcal{G}_i, \psi_a)$
of systems of abelian sheaves. Note that $\mathcal{G}_i$ is an injective
sheaf of abelian groups on $X_i$, see
Lemma \ref{lemma-pushforward-injective-flat} and
Homology, Lemma \ref{homology-lemma-product-injectives}.
This finishes the construction.

\medskip\noindent
Arguing exactly as in the proof of
Lemma \ref{lemma-quasi-separated-cohomology-colimit}
we see that it suffices to prove that
$H^p(X, \colim f_i^{-1}\mathcal{G}_i) = 0$ for $p > 0$.

\medskip\noindent
Set $\mathcal{G} = \colim f_i^{-1}\mathcal{G}_i$.
To show vanishing of cohomology of $\mathcal{G}$ on every quasi-compact
open of $X$, it suffices to show that the {\v C}ech cohomology of
$\mathcal{G}$ for any open covering $\mathcal{U}$ of a quasi-compact open of
$X$ by finitely many quasi-compact opens is zero, see
Lemma \ref{lemma-cech-vanish-basis}.
Such a covering is the inverse by $p_i$ of such a covering $\mathcal{U}_i$
on the space $X_i$ for some $i$ by
Topology, Lemma \ref{topology-lemma-descend-opens}. We have
$$
\check{\mathcal{C}}^\bullet(\mathcal{U}, \mathcal{G}) =
\colim_{a : j \to i}
\check{\mathcal{C}}^\bullet(f_a^{-1}(\mathcal{U}_i), \mathcal{G}_j)
$$
by the case $p = 0$. The right hand side is a filtered colimit of
complexes each of which is acyclic in positive degrees by
Lemma \ref{lemma-injective-trivial-cech}. Thus we conclude by
Algebra, Lemma \ref{algebra-lemma-directed-colimit-exact}.
\end{proof}

\end{document}
